\cleardoublepage
\chapter*{Advertencia}
\thispagestyle{empty}
\label{I.avertissement}

Cada una de las exposiciones redactadas recoge el contenido de varias exposiciones orales consecutivas. No ha parecido útil precisar sus fechas.

La exposición~VII, a la que se remite en diversas ocasiones a lo largo de la exposición~\Ref{VIII}, no fue redactada por el conferenciante, quien, en las conferencias orales, se había limitado a esbozar el lenguaje del descenso en las categorías generales, situándose desde un punto de vista estrictamente utilitario y sin entrar en las dificultades lógicas que plantea este lenguaje. Resultó que una exposición correcta de este lenguaje rebasaría los límites de las presentes notas, aunque solo fuera por su extensión. Para una exposición formal de la teoría del descenso, remito a un artículo en preparación de Jean \textsc{Giraud}. Mientras aparece\footnote{Publicado entretanto: J\ptbl\ \textsc{Giraud}, \emph{\ifthenelse{\boolean{orig}}{Méthodes}{Méthode} de la descente}, Mémoire n.º~2 de la Société mathématique de France, 1964.\label{footnotegiraud}}, pienso que un lector atento no tendrá dificultad en suplir por sus propios medios las referencias fantasma de la exposición~VIII.

Tampoco fueron redactadas otras exposiciones orales, posteriores a la exposición~\Ref{XI} y a las que se alude en algunos lugares del texto; estaban destinadas a constituir el contenido de una exposición~\Ref{XII} y de una exposición~\Ref{XIII}. Las primeras de estas exposiciones orales retomaban, en el marco de los esquemas y de los espacios analíticos con elementos nilpotentes (tal como se introducen en el Seminario Cartan 1960/61), la construcción del espacio analítico asociado a un preesquema localmente de tipo finito sobre un cuerpo valuado completo~$k$, los teoremas de tipo GAGA en el caso en que $k$ es el cuerpo de los números complejos y la aplicación a la comparación entre el grupo fundamental definido por vía trascendente y el grupo fundamental estudiado en estas notas (compárese A\ptbl\ Grothendieck, \emph{Fondements de la Géométrie Algébrique}, Seminario Bourbaki n.º~190, página~10, diciembre de 1959). Las últimas exposiciones orales esbozaban la generalización de los métodos desarrollados en el texto para estudiar los recubrimientos que admiten ramificación moderada y la estructura del grupo fundamental de una curva completa menos un número finito de puntos (compárese \loccit, n.º~182, página~27, teorema~14). Estas exposiciones no introducen ninguna idea esencialmente nueva, razón por la cual no ha parecido indispensable darles una redacción formal antes de la aparición de los capítulos correspondientes de los \emph{Elementos de geometría algebraica}\footnote{Se incluyen en el presente volumen, en la exp\ptbl\ \Ref{XII} de la Sra.~Raynaud, con una demostración distinta de la demostración original expuesta en el Seminario oral (\cf introducción).}.

En cambio, los teoremas de tipo Lefschetz para el grupo fundamental y el grupo de Picard, tanto desde el punto de vista local como desde el global, fueron objeto de un Seminario separado en 1962, que fue redactado por completo y está a disposición de los usuarios\footnote{\emph{Cohomologie locale des faisceaux cohérents et théorèmes de Lefschetz locaux et globaux} (SGA~2), publicado por North-Holland Publishing Company.}. Señalemos que los resultados desarrollados tanto en el presente Seminario como en el de 1962 se utilizarán de manera esencial para la publicación de varios resultados clave de la cohomología étale de los preesquemas, que serán objeto de un Seminario (dirigido por M\ptbl\ Artin y por mí) en 1963/64, actualmente en preparación\footnote{\emph{Cohomologie étale des schémas} (citado como SGA~4), de próxima aparición en esta misma serie.}.

Las exposiciones \Ref{I} a~\Ref{IV}, de naturaleza esencialmente local y muy elementales, quedarán enteramente incorporadas en el capítulo~IV de los \emph{Elementos de geometría algebraica}, cuya primera parte está en prensa y que se publicará sin duda hacia finales de 1964. No obstante, podrán resultar útiles al lector que desee ponerse al corriente de las propiedades esenciales de los morfismos lisos, étales o planos antes de adentrarse en los arcanos de un tratado sistemático. En cuanto a las demás exposiciones, quedarán incorporadas en el capítulo~VIII\footnote{De hecho, a raíz de una modificación del plan inicialmente previsto para los \emph{Elementos}, el estudio del grupo fundamental queda aplazado allí hasta un capítulo posterior al que acaba de indicarse. Compárese la introducción que precede a la presente «Advertencia».} de los «Elementos», cuya publicación difícilmente podrá contemplarse antes de varios años.

\bigskip\hfill Bures, junio de 1963.

{\def\footnotemark{}
\chapterspace{-4}
\tableofcontents}

\mainmatter
\renewcommand{\chaptername}{Exposición}
\let\thesubsection\oldthesubsection

\chapterspace{-2}
